\documentclass{article}
\usepackage[T1]{fontenc}
\usepackage{lmodern}
\usepackage{graphicx}
\usepackage{amsmath,amssymb}
\usepackage[a4paper,margin=2cm]{geometry}
\usepackage[absolute,overlay]{textpos}
\setlength{\TPHorizModule}{1mm}
\setlength{\TPVertModule}{1mm}
\usepackage[indonesian]{babel}
\usepackage{hyperref}
\setlength{\emergencystretch}{2em}
% unit: Halaman 19

\begin{document}

% === Halaman 19 (python/native) ===

19

Kelebihan Mode ECB

1.

Karena tiap blok plainteks dienkripsi secara independen, maka kita 

tidak perlu mengenkripsi pesan secara sekuensial/linier. 

• Kita dapat mengenkripsi 5 blok pertama, kemudian blok-blok di akhir, dan  

kembali ke blok-blok di tengah dan seterusnya.

• Mode ECB cocok untuk mengenkripsi isi arsip (file) yang diakses secara acak, 

misalnya arsip-arsip basisdata. 

• Jika basisdata dienkripsi dengan mode ECB, maka sembarang record dapat 

dienkripsi secara independen dari record lainnya (dengan asumsi setiap 

record terdiri dari sejumlah blok diskrit yang sama banyaknya).

\end{document}
